% TOP_PROBLEM: {"id": 30007290, "title": "Weil structures on regular tame eigensheaves", "category_id": 1, "category": {"id": 1, "name": "number_theory", "display_name": "Number Theory", "description": "Properties of integers, prime numbers, Diophantine equations.", "slug": "number-theory", "order_index": 1, "created_at": "2026-07-31T15:26:25.670Z"}, "status": "open"}
% FIRST_POSTED: 2026-10-06
% ATTEMPT_STATUS: solved
% OPENAI_CLAIM: {"schema_version":1,"source_repo":"https://github.com/openai/math","source_commit":"adc7f1241b42e322a6451854ab7e4b4c146bf78a","release_date":"2026-10-06","match":"full","resolution":"proved","independently_reviewed":false,"second_pass":"agreed","families":[{"family":"014","title":"Restricted geometric Langlands, global Arthur enhancements, and generic Ramanujan","subject":"Number theory","manuscripts":[{"title":"Frobenius Structures on Tame Hecke Eigensheaves","pdf_path":"preprints/Frobenius-Structures-on-Tame-Hecke-Eigensheaves-October-5-2026/tame-hecke-frobenius.pdf","readme_path":"preprints/Frobenius-Structures-on-Tame-Hecke-Eigensheaves-October-5-2026/README.md","date":"2026-10-05","theorem_ref":"thm:main, eq:regular-monodromy, eq:main-eigenmaps (sections/introduction.tex:36)","statement_sources":[{"path":"preprints/Frobenius-Structures-on-Tame-Hecke-Eigensheaves-October-5-2026/build/sections/introduction.tex","line":36,"label":"thm:main, eq:regular-monodromy, eq:main-eigenmaps"}]}],"lean":null,"reasoning_trace":null}]}
\documentclass[11pt]{article}
\usepackage[margin=25mm]{geometry}
\usepackage{fontspec}
\setmainfont{Latin Modern Roman}
\usepackage{amsmath,amssymb,mathtools}
\usepackage{unicode-math}
\setmathfont{Latin Modern Math}
\usepackage{enumitem}
\usepackage{xurl,hyperref,bookmark}
\hypersetup{colorlinks=true,urlcolor=blue}
\setcounter{secnumdepth}{0}
\setlength{\parindent}{0pt}
\setlength{\parskip}{5pt}
\setlength{\emergencystretch}{3em}
\begin{document}
\section*{Weil structures on regular tame eigensheaves}

\subsection{OpenAI's claim}
OpenAI presents the following claims relevant to this question, with manuscript numbers following paper order in Sources. In manuscript 1 (family 014; 2026-10-05), the asserted result is that after a finite constant-field extension, the one-point tame eigensheaf construction for $\mathrm{SL}_n$, $p>n$, can retain the prescribed arithmetic Frobenius and full multi-leg eigenstructure.

\subsection{Mathematical statement of the claimed result}
\paragraph{Manuscript 1.}
Let $n\ge2$, $\mathbb F_q$ have characteristic $p>n$, and $\ell\ne p$, $E=\overline{\mathbb Q}_{\ell}$. Let $X_0/\mathbb F_q$ be smooth projective geometrically connected of genus at least two, with rational point $x_0$, and $U_0=X_0\setminus\{x_0\}$. A continuous $\rho:\pi_1^{\rm et}(U_0)\to\operatorname{PGL}_n(L)$, $L/\mathbb Q_{\ell}$ finite, is assumed geometrically Zariski dense, trivial on wild inertia, and regularly unipotent on tame inertia:
\[
 \rho(\gamma)=\exp(t_{\ell}(\gamma)N),\qquad N\text{ regular nilpotent}.
\]
The claim gives $m\ge1$ and a Weil sheaf $M_0\ne0$ with $E$-coefficients on $\operatorname{Bun}_{\operatorname{SL}_n,B,x_0}(X_0)_{\mathbb F_{q^m}}$. After geometric pullback it is perverse and locally constructible, its singular support lies in the parabolic nilpotent locus, and
\[
 \operatorname{Hecke}_{I,(V_i)}(M_0)\simeq
 M_0\boxtimes\mathop{\boxtimes}_{i\in I}(V_i)_{\rho_m}
\]
for all finite leg sets and dual-group representations. These are Weil isomorphisms coherent for units, convolution, permutations and all fusion diagonals; $\rho_m$ is the restriction to the extended constant field.

\subsection{Outline of the argument}
For manuscript 1 (family 014; 2026-10-05), an arithmetic parameter is lifted to characteristic zero with its boundary action intact and then glued across a separating node. Equivariant nearby cycles and a compact conjugator preserve Frobenius throughout specialization and descent from the enhanced flag level.

\subsection{Scope relative to this problem}
The result supplies the conclusion in this problem with its displayed hypotheses and quantifiers. The broader statements and applications in the manuscript are kept separate from the particular question above. The supporting result is manuscript 1 (family 014; 2026-10-05).

\subsection{Formal verification}
No Lean formalization is listed at this snapshot for manuscript 1 (family 014; 2026-10-05).

This site has not executed the Lean code. Described statement coverage is not an independent proof check.

\subsection{Status caveat}
The manuscripts were produced by an internal OpenAI model. OpenAI cautions that results without formalization can contain mistakes. This site has not independently checked the proofs; the label records the scope of OpenAI's claim.

\subsection{Sources}
\begin{itemize}
\item OpenAI, \emph{Frobenius Structures on Tame Hecke Eigensheaves}, 2026-10-05; citation key \texttt{OAI:Frobenius-Structures-on-Tame-Hecke-Eigensheaves-October-5-2026}. \href{https://github.com/openai/math/blob/adc7f1241b42e322a6451854ab7e4b4c146bf78a/preprints/Frobenius-Structures-on-Tame-Hecke-Eigensheaves-October-5-2026/tame-hecke-frobenius.pdf}{PDF (GitHub)}; \href{https://github.com/openai/math/raw/adc7f1241b42e322a6451854ab7e4b4c146bf78a/preprints/Frobenius-Structures-on-Tame-Hecke-Eigensheaves-October-5-2026/tame-hecke-frobenius.pdf}{PDF (direct)}; \href{https://github.com/openai/math/blob/adc7f1241b42e322a6451854ab7e4b4c146bf78a/preprints/Frobenius-Structures-on-Tame-Hecke-Eigensheaves-October-5-2026/README.md}{Manuscript page and citation}; \href{https://github.com/openai/math/blob/main/preprints/Frobenius-Structures-on-Tame-Hecke-Eigensheaves-October-5-2026/tame-hecke-frobenius.pdf}{Latest version}; Statement source: \href{https://github.com/openai/math/blob/adc7f1241b42e322a6451854ab7e4b4c146bf78a/preprints/Frobenius-Structures-on-Tame-Hecke-Eigensheaves-October-5-2026/build/sections/introduction.tex\#L36}{thm:main, eq:regular-monodromy, eq:main-eigenmaps}.
\item \href{https://github.com/openai/math/blob/adc7f1241b42e322a6451854ab7e4b4c146bf78a/CONTENTS.md}{OpenAI catalogue entry}.
\end{itemize}
\end{document}
